\thispagestyle{plain}
\phantomsection
\addcontentsline{toc}{chapter}{中文摘要}
\centerline{\zihao{-3}\heiti 中文摘要}

\linespread{1.4}\zihao{-4}\bigskip

本文针对“高精度参数确定问题”建立了基于常微分方程的参数辨识模型，并结合压缩包中的\texttt{PredatorPreyProblem.tex}、相关导数与程序文件，对求解流程进行了系统化重构。研究对象可抽象为带观测噪声的非线性动力系统，目标是在保证轨迹拟合精度的同时提升参数估计稳定性。

文中首先给出状态方程、观测方程与残差目标函数，随后构建高斯牛顿与阻尼牛顿两类迭代方法；在数值求解环节，采用高阶泰勒展开/自适应步长积分策略，提高模型输出对参数变化的可辨识性。针对初值敏感、噪声放大和局部极小值问题，进一步设计了多初值筛选、正则化约束和蒙特卡洛扰动实验。

结果表明：在中低噪声场景下，高斯牛顿法可快速达到较低残差；在噪声偏高或初值偏差较大时，阻尼机制可有效抑制发散并改善参数恢复精度。本文形成了“方程建模--参数优化--鲁棒验证”的完整流程，可为复杂动力系统的高精度参数辨识提供参考。

\bigskip
\noindent{\zihao{4}\heiti 关键词：}
参数确定；常微分方程；高斯牛顿法；阻尼牛顿法；鲁棒性分析
